\documentclass[11pt]{article}
\usepackage[utf8]{inputenc}
\usepackage[T1]{fontenc}
\usepackage{lmodern}
\usepackage[spanish]{babel}
\usepackage{amsmath,amssymb}
\usepackage{graphicx}
\usepackage[left=2cm,right=2cm,top=2cm,bottom=2cm]{geometry}
\usepackage{enumitem}
\usepackage{booktabs}
\usepackage{array}
\usepackage{xcolor}
\usepackage{framed}

\spanishdecimal{.}

\setlist[enumerate,1]{label=\textbf{\arabic*.}, leftmargin=*, itemsep=6pt}
\setlist[enumerate,2]{label=(\alph*), itemsep=4pt}
\setlength{\parindent}{0pt}
\setlength{\parskip}{4pt}
\pagestyle{plain}

\newcommand{\Ein}{E_{\text{in}}}
\newcommand{\Eout}{E_{\text{out}}}
\newcommand{\Hyp}{\mathcal{H}}
\newcommand{\dvc}{d_{\text{VC}}}
\newcommand{\vx}{\mathbf{x}}
\newcommand{\vw}{\mathbf{w}}
\newcommand{\vy}{\mathbf{y}}
\newcommand{\R}{\mathbb{R}}
\newcommand{\T}{^{\top}}
\newcommand{\E}{\mathbb{E}}
\newcommand{\Prob}{\mathbb{P}}
\newcommand{\pts}[1]{\hfill\textbf{(#1~puntos)}}
\newcommand{\linea}[1]{\rule{#1}{0.4pt}}
\newcommand{\ptsinc}[1]{{\unskip\nobreak\hfil\penalty50\hskip1em\hbox{}\nobreak\hfill(#1)\parfillskip=0pt\finalhyphendemerits=0\par}}

\newcommand{\encabezado}[2]{%
\begin{center}
{\large\bfseries Reconocimiento de Patrones}\\[2pt]
Universidad de Sonora \textperiodcentered{} Licenciatura en Ciencias de la Computación\\
Semestre 2027-1 \textperiodcentered{} Prof.\ Julio Waissman Vilanova\\[10pt]
{\Large\bfseries #1}%
\ifdefined\tipodoc\\[4pt]{\large\itshape\tipodoc}\fi
\end{center}
\ifdefined\tipodoc\else
\vspace{6pt}
\textbf{Nombre:} \hrulefill\quad \textbf{Expediente:} \rule{3.5cm}{0.4pt}\par
\vspace{6pt}
\fi
#2
\par\medskip\hrule\medskip
}

\newcommand{\espacio}[1]{\par\vspace{#1}}
\newcommand{\nuevapagina}{\newpage}


\begin{document}

\encabezado{Examen rápido 1 \textperiodcentered{} Aprender de datos y preparar atributos}{%
Tiempo: 20 minutos. Responde en esta hoja. Puedes usar calculadora; no se permiten apuntes.}

\begin{enumerate}

%% ---------------------------------------------------------------
\item \textbf{Del registro al vector.} Para predecir el precio de casas tenemos cinco casas de entrenamiento:

\begin{center}
\begin{tabular}{cccc}
\toprule
Casa & Área (m$^2$) & Cocina & Barrio \\
\midrule
A & 70 & Buena & Norte \\
B & 90 & Regular & Centro \\
C & --- & Excelente & Norte \\
D & 110 & Buena & Sur \\
E & 130 & Mala & Centro \\
\bottomrule
\end{tabular}
\end{center}

El \emph{pipeline} trata cada columna así:
\begin{itemize}[itemsep=0pt]
  \item \emph{Área}: imputar los faltantes con la mediana y después estandarizar, $z = (x - \bar x)/s$, con la desviación estándar $s$ calculada con denominador $n$ (como lo hace \texttt{StandardScaler}).
  \item \emph{Cocina}: codificación ordinal, Mala $=0$, Regular $=1$, Buena $=2$, Excelente $=3$.
  \item \emph{Barrio}: \emph{one-hot} con las categorías en orden alfabético; una categoría que no se vio en el entrenamiento se codifica con puros ceros.
\end{itemize}

\begin{enumerate}
  \item ¿Qué valores aprende el \emph{pipeline} al llamar a \texttt{fit} con estas cinco casas? Escribe cada uno.
  \espacio{2.8cm}

  \item ¿Cuántas columnas tiene la matriz que entrega el \emph{pipeline}?
  \espacio{0.8cm}

  \item Escribe el vector que produce \texttt{transform} para cada una de estas casas nuevas:
  \begin{itemize}[itemsep=0pt]
    \item F: área 150 m$^2$, cocina Excelente, barrio Sur.
    \item G: área desconocida, cocina Regular, barrio Pitic.
  \end{itemize}
  \espacio{2.5cm}

  \item Un compañero propone calcular la mediana, la media y la desviación con las siete casas (las de entrenamiento y las de prueba) «para tener más datos». ¿Qué problema tiene la propuesta?
  \espacio{2.5cm}
\end{enumerate}

\nuevapagina
%% ---------------------------------------------------------------
\item \textbf{¿Información legítima?} Un banco quiere un modelo que, \emph{antes de llamar} a un cliente, estime la probabilidad de que contrate un depósito a plazo, para decidir a quién llamar. Tiene los registros de una campaña anterior.

\begin{enumerate}
  \item Marca cada columna como \textbf{L} (se puede usar) o \textbf{F} (fuga de información) y justifica en una línea.

  \begin{center}
  \renewcommand{\arraystretch}{1.6}
  \begin{tabular}{p{6cm}c p{7.6cm}}
  \toprule
  Columna & L / F & Justificación \\
  \midrule
  \texttt{edad} & & \\
  \texttt{saldo}: saldo promedio anual de la cuenta & & \\
  \texttt{resultado\_previo}: resultado con ese cliente en la campaña \emph{anterior} & & \\
  \texttt{duracion}: duración de la llamada, en segundos & & \\
  \texttt{fecha\_alta}: fecha en que se abrió el depósito (vacía si no lo contrató) & & \\
  \bottomrule
  \end{tabular}
  \end{center}

  \item Este código reporta una exactitud que será optimista. Señala \textbf{dos} problemas (por número de línea) y di cómo corregirlos.
\begin{verbatim}
1  X = pd.get_dummies(datos.drop(columns="suscribe"))
2  X = SimpleImputer(strategy="median").fit_transform(X)
3  X_ent, X_pru, y_ent, y_pru = train_test_split(X, datos["suscribe"])
4  for C in [0.01, 0.1, 1, 10, 100]:
5      modelo = LogisticRegression(C=C).fit(X_ent, y_ent)
6      print(C, modelo.score(X_pru, y_pru))   # reporto el mejor
\end{verbatim}
  \espacio{3.2cm}
\end{enumerate}

\nuevapagina
%% ---------------------------------------------------------------
\item \textbf{¿Aprender o programar?} Recuerda las tres condiciones para que tenga sentido aprender: hay un patrón, no sabemos escribirlo y tenemos datos que se parecen a los casos futuros. Para cada situación, si alguna condición falla di cuál; si tiene sentido aprender, di el tipo de aprendizaje y qué serían $\vx$ y $y$ (si hay $y$).
\begin{enumerate}
  \item Calcular la retención del impuesto sobre la renta de cada empleado en la nómina.
  \espacio{2.2cm}
  \item Pronosticar el volumen que entrará a una presa la próxima semana, con el historial de lluvias y escurrimientos.
  \espacio{2.2cm}
  \item Elegir los números de la lotería del domingo a partir de los resultados de los últimos 20 años.
  \espacio{2.2cm}
  \item Un detector de neumonía entrenado con radiografías de un hospital de Boston, que se usará en clínicas rurales de la sierra de Sonora con otros equipos de rayos X.
  \espacio{2.2cm}
  \item Encontrar grupos de clientes con hábitos de compra parecidos, sin categorías definidas de antemano.
  \espacio{2.2cm}
\end{enumerate}

\end{enumerate}


\end{document}
